% ===== §5 Neural Illustration (comprehensive review) =====
\section{A Neural Illustration: The Mechanisms of the Plain, Shared, and Waiting Happiness}
\label{p19:sec:neuro}

\noindent
This chapter offers a second illustrative register, parallel and co-equal with the dynamical one, and \emph{subordinate in the very same way}: it is a witness, not a ground; the philosophical claims of \xref{p19:sec:generativity} and \xref{p19:sec:modes} stand on their own before it; and by \xref{p19:sub:gen-manifold} it charts another region of the background manifold's shape, never the meaning that unfolds upon it. The review that follows is deliberately broad---it surveys the principal neural systems implicated in the phenomenon---precisely so that the breadth may make vivid, at its end, the single boundary every one of these systems shares: none reaches the conferral of meaning (\xref{p19:sub:gen-hermeneutic}). Comprehensiveness here serves subordination, not its overturning.

\paragraph{Honesty clause (placed first).} Throughout, \emph{[E]} marks a claim \textbf{established} in the literature, \emph{[S]} a claim with \textbf{some support}, and \emph{[H]} a \textbf{hypothesis advanced here}. Per \xref{p19:sub:gen-methodology}, neural mechanism is \emph{illumination, not proof}; it lends an analysable side-image to claims already established philosophically, does not ground them, and is never to be passed off as empirical proof of any proposition about happiness. The mapping from mechanism to phenomenon is, in every row below, an \emph{interpretive} mapping---an act in the symbolic (\xref{p19:sub:gen-hermeneutic})---not a reduction.

\subsection{The systems implicated: a structured review}
\label{p19:sub:neuro-systems}

\paragraph{(i) Reward and value: dopamine beyond ``pleasure.''} Mesencephalic dopamine encodes a reward \emph{prediction error}, not reward itself \parencite{schultz1997}~\emph{[E]}; mesolimbic dopamine tracks the \emph{value of work} and sustains \emph{ramping} activity across a delay toward an expected reward \parencite{hamid2016}~\emph{[E]}. Crucially, the dopaminergic ``wanting'' system is dissociable from the hedonic ``liking'' system, whose neural substrate is a small set of cortical and subcortical \emph{hedonic hotspots} \parencite{berridge2015,kringelbach2009}~\emph{[E]}. This dissociation matters to the phenomenon: the plain happiness is closer to ``liking'' (a quiet hedonic tone) than to ``wanting'' (appetitive drivenness)---a neural correlate of the distinction between the mode of fullness and the mode of lack (\xref{p19:sub:gen-twomodes}).

\paragraph{(ii) Prediction and free energy.} Under the free-energy framework, the cortex minimises prediction error through active inference, with \emph{precision} weighting the gain on prediction errors \parencite{friston2010}~\emph{[E]}. Plainness corresponds to persistently low prediction error---but low error is two-faced: it can signify a well-modelled, untroubled world or an informationally impoverished one (boredom). The framework alone cannot tell these apart~\emph{[S]}; the criterion must come from elsewhere (subsection (vii) and \xref{p19:sub:neuro-fullness}).

\paragraph{(iii) The default mode network and self-referential rest.} The DMN is active at rest and in self-referential processing \parencite{raichle2001,northoff2006}~\emph{[E]}; mind-wandering, frequently DMN-mediated, correlates with lowered subjective happiness \parencite{killingsworth2010}~\emph{[E]}; ruminative, negatively self-referential activity is a marker of dysphoric states \parencite{nolenhoeksema2008}~\emph{[E]}. Yet experienced meditators show \emph{reduced} DMN activity and altered DMN connectivity \parencite{brewer2011}~\emph{[E]}, and contemplative training reshapes self-referential and attentional networks \parencite{tang2015,lutz2008}~\emph{[S]}. Thus ``rest'' is in itself neutral: its value turns on the \emph{character} of DMN activity (ruminative versus non-ruminative, open), not on its mere presence.

\paragraph{(iv) Autonomic regulation and safety.} A calm, low-arousal, parasympathetically dominant tone---in polyvagal terms a ventral-vagal state of safety and social engagement \parencite{porges2007}~\emph{[S]}---is the physiological signature of ``safe presence,'' the bodily ground of a being-together that requires no vigilance.

\paragraph{(v) Affiliation and the calm-and-connect axis.} Oxytocinergic and related affiliative systems support non-appetitive bonding and self-soothing under non-noxious sensory contact \parencite{uvnasmoberg2015}~\emph{[S]}; the neurobiology of attachment describes a synchronised, affiliative bond not reducible to reward-seeking \parencite{feldman2017}~\emph{[E]}. This is the neural side-image of Winnicott's ``being quietly in the other's presence'' (\xref{p19:sub:theories-psych})---intimacy without the activation of drive.

\paragraph{(vi) The social brain and intersubjectivity.} Watching the rain \emph{together} engages systems beyond the single brain: shared attention and mentalising networks, and, at the inter-brain level, \emph{neural synchrony} between interacting partners \parencite{hasson2012,dikker2017}~\emph{[S]}. Second-person neuroscience argues that genuine interaction is not reducible to two spectator brains \parencite{redcay2019}~\emph{[S]}---a neural counterpart to the irreducibility of the ``we'' (\xref{p19:sub:gen-relational}).

\paragraph{(vii) The positive low-arousal state.} Across (ii)--(iv), the recurring problem is to distinguish a \emph{full} low-arousal state from an \emph{empty} one. The resolution proposed here (developed in \xref{p19:sub:neuro-fullness}) is that the difference lies not in the level of arousal or error but in the \emph{neuromodulatory tone} and the \emph{character of resting activity}.

\subsection{Two review tables}
\label{p19:sub:neuro-tables}

Table~\ref{p19:tab:neuro-systems} consolidates the systems against the facets of the phenomenon; Table~\ref{p19:tab:neuro-modes} consolidates the central distinction between the two ``uneventful plainnesses.''

\begin{table}[t]
\caption{Neural systems implicated in the plain, shared, and waiting happiness. Evidence: \emph{[E]} established, \emph{[S]} some support, \emph{[H]} hypothesis advanced here. The mapping to facets is interpretive (\xref{p19:sub:gen-hermeneutic}), not reductive.}
\label{p19:tab:neuro-systems}
\small
\begin{tabularx}{\linewidth}{>{\RaggedRight\arraybackslash}p{2.6cm} >{\RaggedRight\arraybackslash}p{3.5cm} >{\RaggedRight\arraybackslash}X >{\centering\arraybackslash}p{1.0cm}}
\toprule
\textbf{System} & \textbf{Core finding} & \textbf{Facet of the phenomenon} & \textbf{Ev.} \\
\midrule
Dopamine / value & Prediction-error coding; delay-period ramping; value of work & The fullness of waiting; ``not a blank'' & [E] \\
Wanting vs liking & Dissociable appetitive and hedonic systems; hedonic hotspots & Mode of fullness (liking) vs mode of lack (wanting) & [E] \\
Free energy & Low prediction error under good model; precision weighting & Plainness as low error; ambiguity of ``low'' & [E]/[S] \\
Default mode network & Rest/self-reference; mind-wandering lowers happiness; meditators show reduced DMN & Two plainnesses turn on DMN character & [E] \\
Autonomic (polyvagal) & Ventral-vagal safety / social-engagement tone & Bodily ground of safe being-together & [S] \\
Affiliation / oxytocin & Non-appetitive bonding; self-soothing; attachment synchrony & Intimacy without drive (Winnicott) & [E]/[S] \\
Social brain / inter-brain & Shared-attention \& mentalising; neural synchrony; second-person interaction & Irreducibility of the ``we'' & [S] \\
\bottomrule
\end{tabularx}
\end{table}

\begin{table}[t]
\caption{The two ``uneventful plainnesses,'' distinguished not by the level of error but by the character beneath it. The word ``saturation'' is deliberately avoided (cf.\ \xref{p19:sub:gen-twomodes}).}
\label{p19:tab:neuro-modes}
\small
\begin{tabularx}{\linewidth}{>{\RaggedRight\arraybackslash}p{3.2cm} >{\RaggedRight\arraybackslash}X >{\RaggedRight\arraybackslash}X}
\toprule
\textbf{Property} & \textbf{Empty plainness (boredom/dysphoria)} & \textbf{Full plainness (the phenomenon)} \\
\midrule
Prediction error & Low & Low \\
Neuromodulatory tone & Depleted / low & Steady \\
DMN activity & Ruminative, negatively self-referential & Non-ruminative, near open awareness \\
Autonomic tone & Withdrawn or dysregulated & Ventral-vagal, safe \\
Relation to the gap & Mode of lack & Mode of fullness \\
Evidence for the contrast & [E] for components & [H] for the specific signature \\
\bottomrule
\end{tabularx}
\end{table}

\subsection{Part A revisited: the value-signalling of waiting}
\label{p19:sub:neuro-waiting}

The sustained delay-period activity of the value system \parencite{schultz1997,hamid2016}~\emph{[E]} gives the philosophical claim of \xref{p19:sub:modes-waiting}---that the fullness of waiting comes from a continual movement beneath an apparent stillness---a nearly literal mechanistic side-image. Indeed, in a delayed-gratification paradigm, dopaminergic activity in the ventral tegmental area \emph{rises steadily across the waiting period}, and optogenetic activation or silencing of these neurons respectively extends or shortens the willingness to wait \parencite{wang2021delayed}~\emph{[E]}---the waiting interval is not a blank but an actively sustained state. And the dissociation of \emph{voluntary, oriented} waiting from \emph{forced, uncontrolled} waiting~\emph{[S]} is anchored in stressor-controllability research: controllable anticipation recruits ventromedial-prefrontal regulation, whereas uncontrollable stress engages subcortical stress circuitry, and cortisol directly blunts the neural anticipation of reward \parencite{wood2021controllability,kinner2016}. Sweet waiting and fretful empty waiting differ, then, not in degree but in circuit---anchoring the ``mode of fullness versus mode of lack'' distinction (\xref{p19:sub:gen-twomodes}) to a separable neural difference.

\subsection{Part B revisited: the mechanism of two plainnesses}
\label{p19:sub:neuro-fullness}

The distinction between the two plainnesses (Table~\ref{p19:tab:neuro-modes}) falls not on the height of the prediction error but on two properties beneath it---the neuromodulatory tone (depleted vs steady) and the character of DMN activity (ruminative vs non-ruminative, near open awareness) \parencite{friston2010,raichle2001,northoff2006,brewer2011}. Two convergent literatures support the contrast: depressive rumination is associated with heightened DMN engagement during negative self-reference \parencite{hamilton2015,nolenhoeksema2008}~\emph{[E]}, while meditation---open-monitoring and choiceless-awareness practices in particular---consistently \emph{reduces} DMN activity and connectivity \parencite{brewer2011,garrison2015}~\emph{[E]}. The components are thus established~\emph{[E]}; the specific synthesis proposed here---that ``full plainness $=$ low error $+$ steady tone $+$ non-ruminative DMN'' constitutes a distinct neural signature of the full, as against the empty, plainness---is advanced as a \emph{hypothesis}~\emph{[S]}, now with convergent support though not yet directly tested in the relevant low-arousal, shared condition.

\subsection{Synthesis, and the single boundary the breadth reveals}
\label{p19:sub:neuro-synth}

The modes of happiness answer to several \emph{healthy} modes of the generative predictive system: the waiting mode $=$ the value circuit's sustained positive signalling across the delay; the abiding mode $=$ low error with steady tone and a non-ruminative default network. Both are side-images of a predictive, embodied, and \emph{interpersonally coupled} system relating to the symbolic--real gap in the \emph{mode of fullness}: the gap is still there (the system is alive, predicting, coupled), but it does not operate in the form of lack-and-drivenness.

\paragraph{The boundary the review makes vivid (returning meaning to philosophy).} Seven systems, dozens of findings---and not one crosses from ``this circuit is in this state'' to ``this means happiness to the subject.'' Every row of Table~\ref{p19:tab:neuro-systems} required an \emph{interpretive} mapping in its third column; that mapping is an act in the symbolic, not a reading-off from the mechanism. The very comprehensiveness of the review is what exhibits the boundary: however many systems are charted, the meaning is not among them. The neural charts a wide region of the manifold (\xref{p19:sub:gen-manifold}); the meaning unfolds, as ever, in the symbolic (\xref{p19:sub:gen-hermeneutic}). This is the chapter's subordination, shown rather than merely asserted.
